\documentclass[11pt,a4paper]{article}
\usepackage[margin=1in]{geometry}
\usepackage{amsmath,amssymb,booktabs,tabularx,microtype}
\usepackage{hyperref}
\hypersetup{colorlinks=true,linkcolor=blue,urlcolor=blue,
 pdftitle={Expanded Appendix: Scalable Spatial Model for Brain MRI Lesion Mask Data}}
\emergencystretch=3em
\newcommand{\R}{\mathbb R}
\newcommand{\E}{\operatorname{E}}
\newcommand{\Var}{\operatorname{Var}}
\newcommand{\Cov}{\operatorname{Cov}}
\newcommand{\diag}{\operatorname{diag}}
\newcommand{\tr}{\operatorname{tr}}
\newcommand{\vr}{\operatorname{vec}_{\rm row}}
\newcolumntype{T}{>{\raggedright\arraybackslash}X}
\numberwithin{equation}{section}
\setcounter{tocdepth}{2}
\title{Expanded Mathematical Appendix\\
 \large Scalable Spatial Model for Brain MRI Lesion Mask Data\\
 \normalsize Companion to \texttt{cbmr\_paper.pdf}}
\author{}
\date{15 September 2026}
\begin{document}
\maketitle
\begin{abstract}
This self-contained appendix develops the model in the supplied manuscript
by Yifan Yu and Thomas E. Nichols, \emph{Scalable Spatial Model for Brain
MRI Lesion Mask Data}. Its central parameters are spatially varying
covariate-effect maps for binary subject images. We derive the Poisson
and independent negative-binomial working scores and Hessians, the
separable initialization, the Kronecker preconditioner, the subject-centered
Taylor approximation, and model-based and subject-clustered inference.
The supplement's QR sandwich and toy slope-map variance are expanded,
including their coefficient-order and dimension qualifications.
The manuscript's mass-univariate logistic comparator is treated separately.
All intermediate dimensions and vectorization conventions are explicit.
Where a displayed manuscript formula needs qualification or correction,
that distinction is stated rather than inherited silently.
This is a handwritten mathematical review document, not a new claim of
complete symbolic or proof-assistant certification.
\end{abstract}
\noindent\textbf{Scope.}
The multi-group coordinate-based meta-regression model from
\texttt{imag.a.1057.pdf} has a different parameterization and is
developed in the separate \texttt{imag\_a\_1057\_appendix.pdf}.
Its moment-matched aggregate and study-clustered NB likelihoods
are not substituted for the independent-cell NB working likelihood here.
\tableofcontents
\clearpage

\section{Paper correspondence, model setup, and dimensions}

\subsection{A map to the supplied manuscript}
\begin{center}\small
\begin{tabularx}{\textwidth}{lT}
\toprule
Manuscript location & Material developed here\\
\midrule
Section 2.1; (2.1)--(2.4) & Fixed spline or QMC spatial basis\\
Section 2.2; (2.5)--(2.7) & Mass-univariate comparator and spatial-effect design\\
(2.8)--(2.10); S1.2--S1.3 & Poisson/NB working likelihoods, binary specialization,
and distinction from a normalized binary model\\
(2.11); Algorithm 1, Stage 1 & Identifiable separable Poisson initialization\\
(2.12)--(2.13); (S2.36)--(S2.37) & Family-specific scoring and exact versus
approximate Kronecker preconditioning\\
(2.14); (S2.38)--(S2.42) & Subject-centered Taylor expansion, dimensions,
and correct weighted basis product\\
(2.15)--(2.18) & Voxel contrasts, Fisher information, and a matching
subject-level sandwich\\
S1.1; (S1.19)--(S1.23) & Logistic and Firth comparator derivatives\\
S2.2; (S2.43)--(S2.50) & Weighted-design QR and triangular sandwich solves\\
S2.3 & Toy slope-map variance, coefficient permutation, and corrected
subject/spatial moment factors\\
\bottomrule
\end{tabularx}
\end{center}
Equation numbers in the left column refer to the supplied manuscript;
numbers elsewhere in this appendix refer to the present document.
The manuscript sometimes labels a full-model likelihood by $\gamma$;
we use $\beta$ for its full spatial coefficient vector and reserve
$\gamma$ for the separable initialization's global coefficients,
except in explicitly labelled local supplementary conventions.

\subsection{Notation crosswalk and deliberate disambiguations}
The main-text conventions are retained rather than importing the
different conventions of \texttt{imag.a.1057.pdf}.
In particular, here $Z_i$ and $B_j$ are \emph{columns}; the rows
of the two design matrices are $Z_i^\top$ and $B_j^\top$.
\begin{center}\small
\begin{tabularx}{\textwidth}{lT}
\toprule
Manuscript symbol & Meaning and use in this appendix\\
\midrule
$M,N,R,P$ & Numbers of subjects, voxels, covariates, and spatial bases\\
$Z,\ B,\ X$ & Designs of dimensions $M\times R$, $N\times P$,
and $MN\times RP$, with $X=Z\otimes B$\\
$Y,\ \mu,\ \eta$ & Subject-major stacked count, mean, and log-mean
vectors of length $MN$\\
$\widetilde Y,\ \widetilde\beta$ & The paper's $M\times N$ data array
and $R\times P$ coefficient reshape from (2.14)\\
$\varphi_j$ & Main-text (2.5)'s $R$-vector of voxelwise effects;
in the spatial model $\varphi_j=(I_R\otimes B_j^\top)\beta$\\
$\xi,\ \gamma$ & Separable initialization coefficients of dimensions
$P$ and $R$, not the full $RP$-vector $\beta$\\
$\mu^B,\ \mu^Z,\ W_B,\ W_Z$ & Separable spatial/study factors
and their diagonal Poisson weight matrices in (2.13)\\
$\widetilde Z,\ \bar\eta,\ \widetilde B$ & Centered covariates,
subject-average log mean, and singly weighted basis from (2.14)\\
$\alpha,\ I(\beta)$ & Global NB dispersion and regression Fisher information\\
$C,\ S,\ W_j$ & Contrast $C\in\R^{S\times R}$, number of tested
constraints, and the \emph{signed} statistic in (2.16)\\
$\bar\eta_C,\ \bar\mu_C$ & Contrast maps $(C\otimes B)\beta$
and their componentwise exponentials, not subject averages\\
\bottomrule
\end{tabularx}
\end{center}
New helpers are explicitly introduced: $\operatorname{vec}_{\rm row}$
fixes the paper's vectorization order; $\widetilde\eta$ and
$\widetilde\mu$ are $M\times N$ arrays; $H=-\nabla^2\ell$ is observed
information; $P_0$ is a fixed full-model preconditioner; and
$Q_j$ names a squared Wald statistic, leaving $W_j$ signed.
We use the auxiliary NB shape $\nu=1/\alpha$ so that $r$ remains
the covariate index.

Several manuscript symbols have different local meanings.
Section 2.1 uses $M$ for the number of Fourier features; we call
that number $P_f$ to distinguish it from subjects.
Its QMC coefficients, also called $\gamma$ there, are not the
Stage 1 global coefficients.
Supplement S1.1 calls its $M\times R$ logistic design $X$ and
its voxelwise coefficient $\beta_j$; our comparator uses the
main-text names $Z$ and $\varphi_j$, with this mapping stated again
in Section~\ref{lesion-logistic}.
Supplement S2.2 reuses $B,M,R,S,Z$ for sandwich and QR quantities;
we distinguish those from the spatial basis, dimensions, contrasts,
and study design below.
Supplement S2.3 reverses the Kronecker order and reuses $\alpha$
as an intercept; its local convention is translated explicitly in
Section~\ref{lesion-toy-variance}.

\subsection{Data and stacking convention}
There are $M$ independent subjects and $N$ voxels.
Subject $i$ supplies a binary image
$Y_i=(Y_{i1},\ldots,Y_{iN})^\top$, $Y_{ij}\in\{0,1\}$.
Independence \emph{within} an image is a working-likelihood assumption,
not a consequence of smoothing its mean.
Let $\widetilde Y\in\R^{M\times N}$ store the images by row.
Stack subjects slowest:
$Y=(Y_1^\top,\ldots,Y_M^\top)^\top=\vr(\widetilde Y)$.
For any matrix, $\vr$ stacks rows in order, equivalently applying
ordinary column-vectorization to its transpose.

The fixed subject design is $Z\in\R^{M\times R}$ and the fixed
spatial basis is $B\in\R^{N\times P}$.
Rows are $Z_i^\top$ and $B_j^\top$; $Z_i\in\R^R$ and $B_j\in\R^P$
are columns. Stack the $R$ spatial maps as
\[
 \beta=(\beta_1^\top,\ldots,\beta_R^\top)^\top\in\R^{RP},
 \quad
 \widetilde\beta=\begin{pmatrix}\beta_1^\top\\ \vdots\\\beta_R^\top\end{pmatrix},
 \quad \beta=\vr(\widetilde\beta).
\]
The model corresponding to manuscript (2.6)--(2.7) is
\begin{align}
 \eta_{ij}=\log\mu_{ij}
 &=\sum_{r=1}^R Z_{ir}B_j^\top\beta_r
   =(Z_i^\top\otimes B_j^\top)\beta,\\
 \widetilde\eta&=Z\widetilde\beta B^\top\in\R^{M\times N},\\
 \eta&=(Z\otimes B)\beta=X\beta,\qquad X\in\R^{MN\times RP}.
 \label{lesion-design}
\end{align}
Here $\widetilde\eta$ is the \emph{array of log means}, not a Hessian.
The study design is $X_i=Z_i^\top\otimes B\in\R^{N\times RP}$.
Let $x_{ij}=Z_i\otimes B_j$ be a cell design column.
The number of regression coefficients is $d=RP$.

The coefficient of covariate $Z_{ir}$ at voxel $j$ is
$\varphi_{jr}=B_j^\top\beta_r$, component $r$ of the paper's
$\varphi_j$. Thus $\partial\eta_{ij}/\partial Z_{ir}=\varphi_{jr}$.
An intercept column of $Z$ supplies a baseline map, not an additional
constant global covariate effect. A common smooth basis reduces the
number of coefficients but does not itself specify residual spatial
covariance between $Y_{ij}$ and $Y_{ik}$.

\subsection{Derive the coefficient order rather than assuming a vec convention}
The cell $(i,j)$ has stacked row index $(i-1)N+j$, and the
coefficient $(r,p)$ has index $(r-1)P+p$.
By the definition of a Kronecker product,
\[
 (Z\otimes B)_{(i-1)N+j,\,(r-1)P+p}=Z_{ir}B_{jp}.
\]
Multiplication by $\beta=\vr(\widetilde\beta)$ therefore gives
$\sum_{r,p}Z_{ir}\widetilde\beta_{rp}B_{jp}$, which is exactly
$(Z\widetilde\beta B^\top)_{ij}$.
This proves both (2.7)'s subject-major stacking and the paper's
$R\times P$ reshape convention.

More generally, if $A$ is $R\times R$, $K$ is $P\times P$,
and $D$ is $R\times P$, then
\[
 (ADK)_{rp}=\sum_{q,l}A_{rq}D_{ql}K_{lp}.
\]
The coefficient multiplying $D_{ql}$ is
$(A\otimes K^\top)_{(r-1)P+p,\,(q-1)P+l}$, proving
\[
 \vr(ADK)=(A\otimes K^\top)\vr(D).
\]
Ordinary column-vectorization of that same $R\times P$ array
would interchange the coefficient order. It is valid only with
the corresponding permutation of the design and covariance.

\subsection{Spline and QMC bases are fixed design choices}
For coordinates $s_j=(x_j,y_j,z_j)^\top$ a tensor spline column
is indexed by $(a,b,c)$:
\[
 B_{j,(a,b,c)}
 =B_a^{(x)}(x_j)B_b^{(y)}(y_j)B_c^{(z)}(z_j).
\]
Then $B_j^\top\beta_r$ is the triple sum of these products times
the corresponding map coefficients, recovering (2.1)--(2.2).
This uses the paper's uppercase univariate basis functions.
Its generic voxel index $i$ in Section 2.1 is changed to $j$ here,
so that $i$ consistently indexes subjects in the regression.
Its generic spline coefficient $\omega_{abc}$ becomes the
corresponding component of $\beta_r$ for covariate map $r$.
On a full rectangular grid the full basis can itself be built by
Kronecker products. Masking voxels or dropping weakly supported
columns does not preserve every such full-grid factorization.
This spatial tensor construction is distinct from $X=Z\otimes B$.

For $P_f$ fixed Fourier features, the manuscript's alternative has
\[
 B_{jk}=\sqrt{2/P_f}\cos(\omega_k^\top s_j+b_k),\quad k=1,\ldots,P_f.
\]
Thus $B_j=\varphi_{\rm QMC}(s_j)$ before any additional intercept
column is appended. The same feature vector in (2.3) is multiplied
by a generic coefficient vector; the spatial GLM supplies one such
vector $\beta_r$ for each covariate.
The frequencies and phases are generated before regression and held
fixed; an added intercept, if independent, increases $P$ by one.
For clarity $P_f$ counts features, avoiding the manuscript Section 2.1
reuse of $M$, which here always counts subjects.
Conditioning on this basis, every regression derivative below is
identical for splines and Fourier features.
No differentiation through a chosen knot grid or QMC sequence is involved.

An intercept should not be appended blindly: a complete partition-of-unity
spline basis already spans the constant function.
If $Bc_B=\mathbf1_N$, adding another constant column makes the basis
rank deficient unless a redundant direction is removed.
For full-column-rank $Z,B$,
$\operatorname{rank}(Z\otimes B)=\operatorname{rank}(Z)\operatorname{rank}(B)=RP$.
Having a constant direction in both factors does not itself duplicate
the intercept in the \emph{tensor-product} design.
The additive initialization has a different intercept problem,
derived in Section~\ref{lesion-initialization}.

\section{Binary observations: count working models versus normalized binary models}
\label{lesion-binary}

\subsection{What conditioning a Poisson count on zero or one changes}
For $Y\sim\operatorname{Poisson}(\lambda)$,
$P(Y=0)=e^{-\lambda}$ and $P(Y=1)=\lambda e^{-\lambda}$.
After explicitly conditioning on $Y\in\{0,1\}$,
\[
 P(Y=1\mid Y\le1)=\frac{\lambda}{1+\lambda},\qquad
 P(Y=0\mid Y\le1)=\frac1{1+\lambda}.
\]
With $\lambda=e^\eta$, the conditional probability is logistic:
$p=e^\eta/(1+e^\eta)$. This is the exact conditioning calculation
underlying manuscript S1.2.
But observing only zeros and ones is not itself an instruction to
condition the count likelihood.
For an observed $y\in\{0,1\}$ the two log-likelihoods are
\begin{align}
 \ell_{\rm count}(\eta)&=y\eta-e^\eta,\\
 \ell_{\rm conditional}(\eta)
 &=\ell_{\rm count}(\eta)
     -\log P_\eta(Y\le1)\\
 &=y\eta-e^\eta-\{-e^\eta+\log(1+e^\eta)\}\\
 &=y\eta-\log(1+e^\eta).
 \label{lesion-conditioned-poisson}
\end{align}
Their difference, $e^\eta-\log(1+e^\eta)$, is parameter-dependent.
Consequently their scores and Hessians differ.
The manuscript's (2.8) uses the unconditioned count objective;
it is not made identical to logistic regression by restricting
the observed dataset to zero and one.

For small $\lambda$,
$\lambda/(1+\lambda)=\lambda-\lambda^2+O(\lambda^3)$ and
$\log(1+\lambda)=\lambda-\lambda^2/2+O(\lambda^3)$.
This gives a low-rate approximation, not equality of the fitted objectives.
Under the unconditioned count model, $\mu=e^\eta$ is an intensity
and is not constrained below one. If it is interpreted as a
Bernoulli mean, fitted values must satisfy that additional scientific
requirement; clipping them after fitting changes neither the fitted
likelihood nor its derivative correctly.

\subsection{The same issue for NB evaluated on binary data}
For an NB count with mean $\mu$ and dispersion $\alpha>0$,
\[
 P(Y=0)=(1+\alpha\mu)^{-1/\alpha},\quad
 \frac{P(Y=1)}{P(Y=0)}=\frac{\mu}{1+\alpha\mu}.
\]
Normalizing these two masses would give
\begin{equation}
 p_{\rm conditional}
 =\frac{\mu}{1+(1+\alpha)\mu},\qquad
 \operatorname{logit}(p_{\rm conditional})
 =\log\mu-\log(1+\alpha\mu).
\end{equation}
That is a different binary model, with a nonlinear predictor in
$\eta=\log\mu$. Its normalization must not be silently added to the
NB working likelihood in manuscript (2.10).
The following NB derivatives refer to that \emph{unconditioned count}
working likelihood evaluated on binary observations.

If the true conditional data are Bernoulli with mean $p$,
$\Var(Y\mid Z)=p(1-p)<p$ for $0<p<1$.
Thus $\mu+\alpha\mu^2$ with $\alpha>0$ cannot be the literal
conditional variance of a binary response with the same mean.
Pooled heterogeneity or within-subject spatial dependence can still
make a working independent model inadequate, but these are not
positive Bernoulli conditional dispersion parameters.
The working-model covariance and a subject-robust covariance therefore
need separate interpretations.

\section{A common chain rule for spatial coefficient maps}
\label{lesion-chain}

\subsection{Differentiate the predictor before the likelihood}
For coordinate $a=(r-1)P+p$,
\[
 \frac{\partial\eta_{ij}}{\partial\beta_a}=Z_{ir}B_{jp}=x_{ij,a},
 \qquad
 \frac{\partial^2\eta_{ij}}{\partial\beta_a\partial\beta_b}=0.
\]
Because $\mu_{ij}=e^{\eta_{ij}}$,
\[
 \partial_{\beta_a}\mu_{ij}=\mu_{ij}x_{ij,a},\quad
 \partial^2_{\beta_a\beta_b}\mu_{ij}=\mu_{ij}x_{ij,a}x_{ij,b}.
\]
For a scalar likelihood $f(\mu)$, this means
\[
 f_{\eta}=\mu f_\mu,\qquad
 f_{\eta\eta}=\mu^2f_{\mu\mu}+\mu f_\mu.
\]
The final term is essential: the log predictor is affine, but
the mean itself is not.

\subsection{Score, log-likelihood Hessian, and observed information}
Write the working log-likelihood as
$\ell(\beta)=\sum_{i,j}f_{ij}(\eta_{ij})$.
Set $s_{ij}=f'_{ij}$ and $w_{ij}=-f''_{ij}$.
Differentiating first and second gives
\begin{align}
 U_a=\partial_{\beta_a}\ell
 &=\sum_{i,j}s_{ij}x_{ij,a},\\
 \partial_{\beta_b}U_a
 &=\sum_{i,j}\{f''_{ij}x_{ij,b}x_{ij,a}
                +s_{ij}\partial_{\beta_b}x_{ij,a}\}\\
 &=-\sum_{i,j}w_{ij}x_{ij,a}x_{ij,b}.
\end{align}
Define $H=-\nabla_\beta^2\ell$ and $I=\E_\beta H$.
Then
\begin{align}
 U&=\sum_i X_i^\top s_i,\\
 H&=\sum_iX_i^\top\diag(w_i)X_i,\\
 H_{\beta_r\beta_{r'}}
 &=\sum_iZ_{ir}Z_{ir'}B^\top\diag(w_i)B,\\
 H_{(r,p),(r',p')}
 &=\sum_i\sum_jZ_{ir}Z_{ir'}B_{jp}B_{jp'}w_{ij}.
 \label{lesion-hessian-entry}
\end{align}
The actual \emph{log-likelihood Hessian} is the negative of every
displayed information block. The scalar weights determine the
observation family; the spatial projection is unchanged.

For implementation, compute $\widetilde B_{jp}=w_{ij}B_{jp}$
and $B^\top\widetilde B$ rather than constructing an $N\times N$
diagonal matrix. Assemble studies or batches without building $X$.
A dense weighted basis product costs $O(NP^2)$ and uses an $N\times P$
workspace. Dense $RP\times RP$ information still costs $O((RP)^2)$
storage; avoiding the large design does not remove that cost.

\subsection{An exact matrix-free Hessian--vector product}
For a direction $v\in\R^{RP}$, let
$\widetilde v\in\R^{R\times P}$ satisfy $v=\vr(\widetilde v)$.
The directional change in the log mean has array
\[
 \widetilde q=Z\widetilde v B^\top,\qquad
 \widetilde q_{ij}=x_{ij}^\top v.
\]
Multiply \eqref{lesion-hessian-entry}'s full sum by $v$:
\begin{align}
 Hv&=\sum_{i,j}w_{ij}x_{ij}(x_{ij}^\top v)\\
   &=\vr\!\left[
       Z^\top\{\widetilde w\mathbin{\odot}
                     (Z\widetilde v B^\top)\}B\right],
 \qquad \widetilde w=(w_{ij})_{i,j}.
 \label{lesion-hvp}
\end{align}
The symbol $\odot$ is elementwise multiplication; $\widetilde w$
is an array, whereas the paper's $W$ is the stacked diagonal
matrix. This identity is exact for observed Poisson or NB
regression curvature at fixed dispersion. Substituting Fisher
weights gives $Iv$ instead.
Its quadratic form is
$v^\top Hv=\sum_{ij}w_{ij}(x_{ij}^\top v)^2$, displaying symmetry
and positive semidefiniteness without forming any dense Hessian.

Subject/voxel batches can accumulate the projected products without
storing $X$, $H$, or all of $\widetilde q$.
With the indicated multiplication order, one dense product costs
$O(MRP+MNP)$, aside from obtaining the current weights.
For identified SPD information, an iterative linear solver with
these products and the separable preconditioner can compute a
direction or selected covariance column without dense information
storage. This changes the numerical solve, not the likelihood.
Residual convergence and information rank still matter; the
preconditioner alone is not the final covariance.

\section{Poisson working likelihood: scalar and spatial Hessian derivations}

\subsection{From the mass function to the second derivative}
The mass function and log-likelihood are
\[
 P(Y=y)=\frac{e^{-\mu}\mu^y}{y!},\qquad
 f(\eta)=y\eta-e^\eta-\log\Gamma(y+1).
\]
When $y=0$ or $1$, $\log\Gamma(y+1)=0$, giving (2.8).
Differentiate at fixed data:
\begin{align}
 f_\eta&=y-e^\eta=y-\mu,\\
 f_{\eta\eta}&=-e^\eta=-\mu,\\
 U&=\sum_{i,j}(Z_i\otimes B_j)(Y_{ij}-\mu_{ij}),\\
 H&=\sum_i(Z_iZ_i^\top)\otimes
                     \{B^\top\diag(\mu_i)B\}.
 \label{lesion-poisson-h}
\end{align}
No count appears explicitly in $H$ at a fixed parameter value.
$H(\widehat\beta)$ still depends on the data through the fitted coefficient.
The count-model expectation gives $I=H$.

Under the correctly specified independent Poisson model,
\[
 \Var(U)=\sum_{i,j}x_{ij}x_{ij}^\top\Var(Y_{ij})
        =\sum_{i,j}\mu_{ij}x_{ij}x_{ij}^\top=I.
\]
Under actual binary or spatially correlated data, the algebraic
Hessian is unchanged but this information equality need not hold.

\subsection{Exact score as a small projected matrix}
Let $\widetilde\mu_{ij}=\mu_{ij}$ and $R_Y=\widetilde Y-\widetilde\mu$.
For the $(r,p)$ coefficient the score is
$\sum_i Z_{ir}\sum_j(R_Y)_{ij}B_{jp}$.
This is entry $(r,p)$ of $Z^\top R_YB$, so
\begin{equation}\label{lesion-score-matrix}
 U=\vr\{Z^\top(\widetilde Y-\widetilde\mu)B\}.
\end{equation}
In particular $X^\top Y=\vr(Z^\top\widetilde YB)$.
This equality fixes the vectorization convention left implicit in some
of the manuscript's formulas; using column-wise $\operatorname{vec}$
of the $R\times P$ matrix would order the coefficients differently.
The count projection can be precomputed for a fixed design, while the
mean projection must be updated at each exact score evaluation.

\section{Independent NB working likelihood: binary simplification and derivatives}
\label{lesion-nb}

\subsection{General NB2 likelihood and exact binary cancellation}
Let $\nu=\alpha^{-1}$, with the manuscript's single dispersion
$\alpha>0$ shared over subjects and voxels. The NB2 mass is
\[
 P(Y=y)=\frac{\Gamma(y+\nu)}{\Gamma(\nu)\Gamma(y+1)}
      \left(\frac{\nu}{\nu+\mu}\right)^\nu
      \left(\frac\mu{\nu+\mu}\right)^y.
\]
Its mean and variance under that count model are $\mu$ and
$\mu+\alpha\mu^2$.
The scalar log-likelihood is
\begin{align}
 f={}&\log\Gamma(y+\nu)-\log\Gamma(\nu)-\log\Gamma(y+1)
       +\nu\log\nu+y\eta-(y+\nu)\log(\nu+e^\eta).
\end{align}
For binary $y$,
$\log\Gamma(y+\nu)-\log\Gamma(\nu)=y\log\nu$.
Also $\log\Gamma(y+1)=0$ and
$\log(\nu+\mu)=\log\nu+\log(1+\alpha\mu)$.
Substituting each term,
\begin{align}
 f&=y\log\nu+\nu\log\nu+y\eta
          -(y+\nu)\{\log\nu+\log(1+\alpha\mu)\}\\
  &=y\eta-(y+\alpha^{-1})\log(1+\alpha\mu)\\
  &=y\log\left(\frac\mu{1+\alpha\mu}\right)
                  -\alpha^{-1}\log(1+\alpha\mu).
 \label{lesion-binary-nb}
\end{align}
This derives manuscript (2.10) and (S1.34), including the cancellations.
It is an exact simplification of the NB \emph{count mass} at
binary observations, not a change in the normalization of that mass.

\subsection{First derivative at fixed original dispersion}
Since $\partial_\eta\mu=\mu$,
\begin{align}
 f_\eta
 &=y-(y+\alpha^{-1})\frac{\alpha\mu}{1+\alpha\mu}\\
 &=\frac{y(1+\alpha\mu)-(1+\alpha y)\mu}{1+\alpha\mu}\\
 &=\frac{y-\mu}{1+\alpha\mu}.
 \label{lesion-nb-score}
\end{align}
The denominator is present even for binary $y$ and even though
the link is logarithmic.

\subsection{Quotient rule for observed curvature}
Differentiate numerator and denominator separately:
\begin{align}
 f_{\eta\eta}
 &=\frac{(-\mu)(1+\alpha\mu)-(y-\mu)\alpha\mu}
           {(1+\alpha\mu)^2}\\
 &=-\frac{\mu+\alpha\mu^2+\alpha\mu y-\alpha\mu^2}
                {(1+\alpha\mu)^2}\\
 &=-\frac{\mu(1+\alpha y)}{(1+\alpha\mu)^2}.
\end{align}
Thus
\begin{align}
 w_{ij}^{\rm obs}
 &=\frac{\mu_{ij}(1+\alpha Y_{ij})}{(1+\alpha\mu_{ij})^2},\\
 H_{\beta_r\beta_{r'}}
 &=\sum_i Z_{ir}Z_{ir'} B^\top\diag(w_i^{\rm obs})B,\\
 U_{\beta_r}
 &=\sum_iZ_{ir}B^\top
          \left\{\frac{Y_i-\mu_i}{1+\alpha\mu_i}\right\}.
 \label{lesion-nb-full}
\end{align}
Vector division is elementwise. These are independent-cell NB
derivatives; there is no aggregation across subjects and no shared
gamma multiplier inducing a rank-one voxel correction.
For fixed positive $\alpha$, $w^{\rm obs}>0$.
Consequently the negative log-likelihood is convex in the affine
regression coefficients and strictly convex when $X$ has full column
rank. Existence of a finite optimum and joint convexity with dispersion
are separate questions.

\subsection{Expected information and what it assumes}
Taking an expectation under the NB working distribution gives
\begin{align}
 w^{\rm F}
 &=\frac{\mu(1+\alpha\E Y)}{(1+\alpha\mu)^2}
   =\frac{\mu}{1+\alpha\mu},\\
 I&=\sum_i X_i^\top\diag\left(\frac{\mu_i}{1+\alpha\mu_i}\right)X_i.
 \label{lesion-nb-fisher}
\end{align}
These are manuscript (2.17)'s weights. Independently,
\[
 \Var_{\rm NB}(f_\eta)
 =\frac{\mu+\alpha\mu^2}{(1+\alpha\mu)^2}
 =\frac{\mu}{1+\alpha\mu}.
\]
The expectation of the observed weight uses only the mean, so it
also gives the expected sensitivity of this working score under
another distribution with the same mean.
The variance identity does \emph{not} then follow automatically.
For a truly independent Bernoulli cell with mean $\mu<1$,
\[
 \Var_{\rm Bernoulli}(f_\eta)
 =\frac{\mu(1-\mu)}{(1+\alpha\mu)^2},
\]
and a whole subject may have additional covariance between its cells.
This is precisely why sensitivity and score variability must be
separated in a robust calculation.

\subsection{Dispersion derivatives, transformations, and the boundary}
Using the general NB likelihood, define
$\psi(t)=d\log\Gamma(t)/dt$ and $\psi_1(t)=d\psi(t)/dt$.
At fixed $\mu$,
\begin{align}
 f_\nu&=\psi(y+\nu)-\psi(\nu)+\log\nu+1-\log(\nu+\mu)
                                       -\frac{y+\nu}{\nu+\mu},\\
 f_{\nu\nu}&=\psi_1(y+\nu)-\psi_1(\nu)+\frac1\nu-\frac1{\nu+\mu}
                                       -\frac{\mu-y}{(\nu+\mu)^2},\\
 f_{\eta\nu}
 &=\partial_\nu\left\{\frac{\nu(y-\mu)}{\nu+\mu}\right\}
   =\frac{\mu(y-\mu)}{(\nu+\mu)^2}.
 \label{lesion-nuisance-r}
\end{align}
For the binary-specialized likelihood \eqref{lesion-binary-nb},
one can also avoid the gamma functions entirely. Put $h=1+\alpha\mu$:
\begin{align}
 f_\alpha
 &=\alpha^{-2}\log h-(y+\alpha^{-1})\mu/h,\\
 f_{\alpha\alpha}
 &=-2\alpha^{-3}\log h+2\alpha^{-2}\mu/h
                       +(y+\alpha^{-1})\mu^2/h^2,\\
 f_{\eta\alpha}&=-\mu(y-\mu)/h^2.
 \label{lesion-nuisance-alpha}
\end{align}
For example the two $\alpha^{-2}\mu/h$ terms in the second
derivative come from differentiating $\log h$ in the first term
and $\alpha^{-1}$ in the second.
These formulas agree with $\nu=1/\alpha$ and the full chain rule.
Assemble $H_{\beta\alpha}=-\sum_{i,j}x_{ij}f_{\eta\alpha}$
and $H_{\alpha\alpha}=-\sum_{i,j}f_{\alpha\alpha}$.

For $\delta=\log\alpha$, differentiation gives
\[
 f_\delta=\alpha f_\alpha,\quad
 f_{\delta\delta}=\alpha^2f_{\alpha\alpha}+\alpha f_\alpha,
 \quad H_{\delta\delta}
 =\alpha^2H_{\alpha\alpha}-\alpha U_\alpha.
\]
The extra score term is not zero away from an interior stationary point.
Holding $\widehat\alpha$ fixed while calculating
\eqref{lesion-nb-full} does not propagate dispersion uncertainty.
The Poisson limit $\alpha\downarrow0$ recovers $s=y-\mu$ and
$w^{\rm obs}=w^{\rm F}=\mu$, but $\alpha=0$ is a boundary,
not an interior value of $\delta$. Joint Wald inference at that
boundary is not justified by the preceding interior derivatives.

\section{The manuscript's separable initialization, not its final model}
\label{lesion-initialization}

\subsection{Additive predictor and identification}
Stage 1 of manuscript Algorithm 1 uses a Poisson fit with
\[
 \eta_{ij}^{(0)}=B_j^\top\xi+Z_i^\top\gamma,\qquad
 \mu_j^B=e^{B_j^\top\xi},\quad \mu_i^Z=e^{Z_i^\top\gamma},\quad
 \mu_{ij}^{(0)}=\mu_i^Z\mu_j^B.
\]
These are the manuscript's $\mu^B$ and $\mu^Z$, not the
$\mu^X$ factors used in the other paper.
The dimensionally explicit version of (2.11) is
\[
 \eta^{(0)}=(\mathbf1_M\otimes B)\xi
                      +(Z\otimes\mathbf1_N)\gamma.
\]
The unstacked components $\eta_B=B\xi$ and $\eta_Z=Z\gamma$
have lengths $N$ and $M$ and must be broadcast this way before
they can be added.
This has $P+R$ nominal coefficients, not $PR$.
Here $\gamma$ is globally constant across voxels and is only the
initialization parameterization; the final $\beta_r$ are spatial maps.

If $Bc_B=\mathbf1_N$ and $Zc_Z=\mathbf1_M$, the transformation
$\xi\mapsto\xi+tc_B$, $\gamma\mapsto\gamma-tc_Z$ leaves every
$\eta_{ij}^{(0)}$ unchanged.
Therefore an unconstrained additive fit with both intercepts is
unidentified. Remove the global intercept for Stage 1 or impose an
explicit equivalent constraint.
Centering the nonconstant covariates does not remove a separately
included global intercept.

\subsection{Exact compression into two count marginals}
Introduce auxiliary count margins
$Y_B=\widetilde Y^\top\mathbf1_M$ and
$Y_Z=\widetilde Y\mathbf1_N$, of lengths $N$ and $M$.
Thus $(Y_B)_j=\sum_iY_{ij}$ and $(Y_Z)_i=\sum_jY_{ij}$.
Define $h_B=\sum_j\mu_j^B$ and $h_Z=\sum_i\mu_i^Z$.
Substituting the additive predictor into the full Poisson objective,
\begin{align}
 \ell_0
 &=\sum_{i,j}Y_{ij}(B_j^\top\xi+Z_i^\top\gamma)
                         -\sum_{i,j}\mu_i^Z\mu_j^B+c(Y)\\
 &=Y_B^\top B\xi+Y_Z^\top Z\gamma-h_Bh_Z+c(Y).
 \label{lesion-initial-ll}
\end{align}
The two count marginals suffice for evaluating this likelihood.
Fitting only $(Y_B)_j\sim\operatorname{Poisson}(h_Z\mu_j^B)$
would drop $\sum_i (Y_Z)_i\log(\mu_i^Z/h_Z)$, which generally depends on
$\gamma$. Exact compression is not the same as discarding the
allocation of counts among subjects.

\subsection{Scores and all initialization Hessian blocks}
Set
\begin{align*}
 W_B&=\diag(\mu^B),& W_Z&=\diag(\mu^Z),\\
 a&=B^\top\mu^B,& t&=Z^\top\mu^Z,\\
 K_B&=B^\top W_BB,& K_Z&=Z^\top W_ZZ.
\end{align*}
Exponential differentiation gives
$\nabla_\xi h_B=a$, $\nabla_\xi^2h_B=K_B$,
$\nabla_\gamma h_Z=t$, and $\nabla_\gamma^2h_Z=K_Z$.
Differentiate \eqref{lesion-initial-ll} once:
\[
 U_\xi=B^\top Y_B-h_Za,\qquad
 U_\gamma=Z^\top Y_Z-h_Bt.
\]
Then differentiate the two scores:
\begin{align}
 -\partial_\xi U_\xi&=h_ZK_B,\qquad
 -\partial_\gamma U_\xi=a t^\top,\\
 -\partial_\gamma U_\gamma&=h_BK_Z,\\
 H_0&=\begin{pmatrix}h_ZK_B&a t^\top\\t a^\top&h_BK_Z\end{pmatrix}.
\end{align}
These formulas are written in nominal $(\xi,\gamma)$ coordinates.
If the global intercept is redundant, impose the chosen constraint
before solving: for $(\xi^\top,\gamma^\top)^\top=J_0\vartheta$
with fixed full-column-rank constraint map $J_0$, use score
$J_0^\top(U_\xi^\top,U_\gamma^\top)^\top$ and information
$J_0^\top H_0J_0$. Deleting a redundant coordinate is one such map.
The constrained information is not the $PR\times PR$ information
of the full spatially varying model at the same means.

\subsection{Embedding a separable fit in the full coefficient maps}
Suppose $Ze_0=\mathbf1_M$ selects the full model's intercept column,
and $Bc_B=\mathbf1_N$.
Pad the identifiable Stage 1 global vector with zero in its
removed intercept coordinate and call it $\widehat\gamma$.
Define
\begin{equation}
 \widetilde\beta^{(0)}=e_0\widehat\xi^\top+\widehat\gamma c_B^\top,\qquad
 \beta^{(0)}=\vr(\widetilde\beta^{(0)}).
 \label{lesion-embedding}
\end{equation}
Multiplication checks the initialization:
\[
 Z\widetilde\beta^{(0)}B^\top
 =\mathbf1_M(B\widehat\xi)^\top
           +(Z\widehat\gamma)\mathbf1_N^\top.
\]
Thus it exactly reproduces the separable log means.
The shorthand $\beta^{(0)}\leftarrow\widehat\xi$ in Algorithm 1
cannot literally assign a $P$-vector to a $PR$-vector.
Equation \eqref{lesion-embedding} supplies a dimensionally complete
embedding that also retains the fitted global effects.
If the required intercept directions are not spanned, an explicit
projection is needed and the embedding is no longer exact.
Stage 1 is Poisson, so it does not itself estimate an NB dispersion;
Stage 2 must specify a positive starting $\alpha$ or a separate
dispersion fit rather than inherit a nonexistent Poisson estimate.

\section{Exact scoring and the Kronecker preconditioner}
\label{lesion-preconditioner}

\subsection{Newton and Fisher directions with the correct family score}
A Taylor expansion of the score equation gives
$0\approx U(\beta)-H(\beta)\Delta$, hence the Newton direction
$\Delta=H^{-1}U$.
Fisher scoring uses $\Delta=I^{-1}U$ instead.
The two family scores are
\[
 U_{\rm P}=X^\top(Y-\mu),\qquad
 U_{\rm NB}=X^\top\diag\{(1+\alpha\mu)^{-1}\}(Y-\mu).
\]
For either family define the IRLS working response
$z^*=\eta+(Y-\mu)/\mu$.
For Poisson $W=\diag(\mu)$ and for NB
$W=\diag\{\mu/(1+\alpha\mu)\}$.
Then
\[
 X^\top W(z^*-\eta)
 =\begin{cases}
 X^\top(Y-\mu),&\text{Poisson},\\
 X^\top\{(Y-\mu)/(1+\alpha\mu)\},&\text{NB}.
 \end{cases}
\]
Solving $X^\top WX\,\Delta=X^\top W(z^*-\eta)$ therefore
reproduces the family-specific Fisher direction.
The numerator in manuscript (2.12) and Algorithm 1 is correct
for Poisson, but must carry the NB residual denominator if it
is intended as NB likelihood Fisher scoring.

\subsection{Deriving the exact factorization at separable Poisson means}
At $\mu_{ij}^{(0)}=\mu_i^Z\mu_j^B$, the stacked diagonal weights satisfy
$W_0=W_Z\otimes W_B=\diag(\mu^Z)\otimes\diag(\mu^B)$.
Using $(A\otimes B)(C\otimes D)=AC\otimes BD$ twice,
\begin{align}
 P_0=X^\top W_0X
 &=(Z^\top\otimes B^\top)
       \{W_Z\otimes W_B\}(Z\otimes B)\\
 &=\{Z^\top W_ZZ\}\otimes\{B^\top W_BB\}\\
 &=K_Z\otimes K_B.
 \label{lesion-kron}
\end{align}
If both factors are nonsingular, multiplication proves
$P_0^{-1}=K_Z^{-1}\otimes K_B^{-1}$.
This is the $PR$-dimensional full-model Poisson information
\emph{evaluated at} the separable means, not $H_0$ above.
When later means cease to be separable, it is a fixed preconditioner,
not the current exact information.

One can solve $P_0d=g$ without forming the Kronecker matrix.
Reshape $g=\operatorname{vec}(G)$ and $d=\operatorname{vec}(D)$
with $G,D\in\R^{P\times R}$, so their column stacking matches
the map-major $\beta$ order.
Then
\[
 (K_Z\otimes K_B)\operatorname{vec}(D)
  =\operatorname{vec}(K_BDK_Z^\top)=\operatorname{vec}(G).
\]
Solve $K_BW=G$, followed by $K_ZD^\top=W^\top$.
Factorization costs $O(P^3+R^3)$, with $O(P^2+R^2)$ factor
storage, rather than a dense $O((PR)^3)$ factorization.

\subsection{Why separable NB means do not give these exact weights}
NB Fisher weights at the same means are
$w_{ij}=\mu_i^Z\mu_j^B/(1+\alpha\mu_i^Z\mu_j^B)$.
For study factors $a,b$ and voxel factors $c,d$, their determinant is
\begin{align}
 &\frac{ac}{1+\alpha ac}\frac{bd}{1+\alpha bd}
   -\frac{ad}{1+\alpha ad}\frac{bc}{1+\alpha bc}\\
 &=\frac{\alpha abcd(a-b)(d-c)}
 {(1+\alpha ac)(1+\alpha ad)(1+\alpha bc)(1+\alpha bd)}.
\end{align}
It is generally nonzero, whereas a separable two-by-two weight grid
has determinant zero. Thus manuscript (2.13)'s Poisson products
are not an exact NB weight factorization.
One may deliberately use $P_0$ as an SPD approximation for NB,
but the NB score and final information must retain their own formulas.

\subsection{What a fixed preconditioner can and cannot guarantee}
For a maximized objective $\ell$ and $P_0\succ0$, the direction
$d=P_0^{-1}U$ satisfies $U^\top d=U^\top P_0^{-1}U>0$
unless $U=0$. This is a local ascent direction.
A full step of length one need not increase $\ell$.
For example, maximize $\ell(b)=-h b^2/2$ using scalar $P_0=p>0$:
\[
 b_{\rm new}=b+p^{-1}(-hb)=(1-h/p)b.
\]
Convergence requires $|1-h/p|<1$, equivalently $0<h/p<2$.
Positive $p$ alone does not imply that inequality.
Line search, damping, or other suitable optimization conditions are
needed for a convergence guarantee. The same caveat applies before
introducing a Taylor approximation to the score.

\section{Subject-centered Taylor approximation: complete derivation}
\label{lesion-taylor}

\subsection{Average log mean and centered predictor}
Define
\[
 \bar z=Z^\top\mathbf1_M/M,\quad
 Q_M=I_M-\mathbf1_M\mathbf1_M^\top/M,\quad
 \widetilde Z=Q_MZ,\quad \bar\eta=B\widetilde\beta^\top\bar z.
\]
Here $\widetilde Z$ has the paper's intended meaning: the centered
$M\times R$ covariate matrix. The centering matrix $Q_M$ alone
is $M\times M$ and is not $\widetilde Z$; the right multiplication
by $Z$ is needed in the definition printed after (2.14).
Then
\[
 Z=\mathbf1_M\bar z^\top+\widetilde Z,\qquad
 \widetilde\eta=\mathbf1_M\bar\eta^\top+\Delta,\qquad
 \Delta=\widetilde Z\widetilde\beta B^\top.
\]
For each voxel, $\sum_i\Delta_{ij}=0$.
Let $a_j=e^{\bar\eta_j}$ and $a=\exp(\bar\eta)$.
Importantly,
$a_j=\exp(M^{-1}\sum_i\eta_{ij})$ is not generally the
subject-average mean $M^{-1}\sum_i e^{\eta_{ij}}$.
Jensen's inequality gives $a_j\le M^{-1}\sum_i\mu_{ij}$.

\subsection{Scalar Taylor expansion with a remainder}
Because $\mu_{ij}=a_j e^{\Delta_{ij}}$, expand the scalar exponential:
\begin{align}
 \mu_{ij}
 &=a_j(1+\Delta_{ij})+\mathcal R_{ij},\\
 \mathcal R_{ij}
 &=\frac12a_j e^{t_{ij}}\Delta_{ij}^2,\qquad
 t_{ij}\text{ lies between }0\text{ and }\Delta_{ij},\\
 0\le\mathcal R_{ij}
 &\le\frac12a_j e^{|\Delta_{ij}|}\Delta_{ij}^2.
 \label{lesion-taylor-remainder}
\end{align}
In matrix form, with $D_a=\diag(a)$,
\begin{equation}
 \widetilde\mu_{\rm lin}
 =\mathbf1_M a^\top+\Delta D_a
 =\mathbf1_M a^\top+\widetilde Z\widetilde\beta B^\top D_a.
 \label{lesion-linear-mean}
\end{equation}
The multiplication by $D_a$ weights columns (voxels), not subject rows.
This is exactly the first-order expansion in manuscript S2.1 before
the final matrix simplification.
It can yield negative approximate means if $\Delta_{ij}<-1$.
It should then not be treated as a valid new Poisson or NB mean
inside a log-likelihood.

\subsection{Projecting the approximate Poisson score}
Start with \eqref{lesion-score-matrix} and substitute
\eqref{lesion-linear-mean}:
\begin{align}
 Z^\top(\widetilde Y-\widetilde\mu_{\rm lin})B
 ={}&Z^\top\widetilde YB
 -(Z^\top\mathbf1_M)(a^\top B)\\
 &-(Z^\top \widetilde Z)\widetilde\beta(B^\top D_aB).
 \label{lesion-taylor-matrix}
\end{align}
Every term has dimension $R\times P$.
For $A\in\R^{R\times R}$ and $K\in\R^{P\times P}$,
\[
 \vr(A\widetilde\beta K)=(A\otimes K^\top)\vr(\widetilde\beta).
\]
Here $K_a=B^\top D_aB$ is symmetric, so row vectorization gives
\begin{equation}\label{lesion-taylor-score}
 U_{\rm lin}
 =\vr(Z^\top\widetilde YB)
 -(Z^\top\mathbf1_M)\otimes(B^\top a)
 -\{(Z^\top \widetilde Z)\otimes K_a\}\beta .
\end{equation}
Also $Z^\top \widetilde Z=\widetilde Z^\top \widetilde Z$ because $Q_M^\top Q_M=Q_M$.
The dimensions are $RP\times1$ for every term.
The current $\bar\eta,a,K_a$ depend on the current $\beta$;
this is a residual approximation at an iterate, not an assertion
that its derivative is just the final Kronecker factor.

\subsection{Differentiate the moving expansion point as well}
To derive that derivative, put
$\bar x_j=\bar z\otimes B_j$.
Then $\bar\eta_j=\bar x_j^\top\beta$ and
$\Delta_{ij}=(x_{ij}-\bar x_j)^\top\beta$.
Since $a_j=e^{\bar\eta_j}$, the product rule gives
\begin{align}
 \nabla_\beta\{a_j(1+\Delta_{ij})\}
 &=a_j(1+\Delta_{ij})\bar x_j+
                         a_j(x_{ij}-\bar x_j)\\
 &=a_j(x_{ij}+\Delta_{ij}\bar x_j).
\end{align}
Therefore the negative Jacobian of the approximate score is
\begin{equation}
 A_{\rm lin}=-\frac{\partial U_{\rm lin}}{\partial\beta^\top}
 =\sum_{i,j}a_jx_{ij}(x_{ij}+\Delta_{ij}\bar x_j)^\top .
 \label{lesion-taylor-jacobian}
\end{equation}
The second term need not be symmetric, so $U_{\rm lin}$ is not
in general the gradient of any scalar twice-differentiable objective.
It should not be assigned a ``Hessian'' by retaining only
$(Z^\top\widetilde Z)\otimes K_a$.

For a direct check, take one voxel, two subject rows $(1,-1)$
and $(1,1)$, and parameters $(\eta_0,b)$.
With $a=e^{\eta_0}$ the approximate score is
\[
 U_{\rm lin}=
 \begin{pmatrix}Y_1+Y_2-2a\\-Y_1+Y_2-2ab\end{pmatrix},
 \qquad
 A_{\rm lin}=\begin{pmatrix}2a&0\\2ab&2a\end{pmatrix}.
\]
For $b\ne0$ the mixed derivatives differ.
If this approximate score defines the final estimator, its
estimating-equation sensitivity is \eqref{lesion-taylor-jacobian},
not automatically the exact likelihood's information.
A likelihood-based final covariance requires an exact-score
refinement or justification that the approximation error is
negligible on the relevant inferential scale.
As with any working estimating equation, a sandwich still requires
a well-defined target and appropriate centered score moments.

\subsection{The weighted basis must appear once, not squared}
Define $\widetilde B=D_aB$, which has dimension $N\times P$.
Then the last factor in \eqref{lesion-taylor-score} is
\[
 K_a=B^\top\widetilde B=B^\top D_aB.
\]
In contrast,
$\widetilde B^\top\widetilde B=B^\top D_a^2B$
would weight by $a_j^2$, not $a_j$.
The product $B^\top\widetilde B$ printed in (2.14), Algorithm 1,
and (S2.42) already has the correct single weight under
$\widetilde B=D_aB$. It should not be reinterpreted as
$\widetilde B^\top\widetilde B$.
An alternative square-root-weighted definition
$\widehat B=D_a^{1/2}B$ would permit
$\widehat B^\top\widehat B=B^\top D_aB$.
Neither weighted basis has dimension $N\times N$ unless $P=N$.

The fixed data projection remains $\vr(Z^\top\widetilde YB)$,
using the \emph{unweighted} $B$; weighting that data basis would
change the score. Also the final $RP\times RP$ Kronecker factor
multiplies the vector $\beta$, not its $R\times P$ reshape
$\widetilde\beta$. The latter belongs inside the matrix product
in \eqref{lesion-taylor-matrix}, followed by row-vectorization.
These distinctions make (2.14)'s data projection and reshaped
coefficient notation consistent with Algorithm 1 and (S2.42);
they do not call for replacing the already correct
$B^\top\widetilde B$ factor.

\subsection{Error in the score and a checkable small example}
The exact and approximate Poisson scores differ by
\[
 U-U_{\rm lin}=-\sum_{i,j}x_{ij}\mathcal R_{ij}.
\]
For any compatible vector norm, triangle inequality and
\eqref{lesion-taylor-remainder} yield
\[
 \|U-U_{\rm lin}\|
 \le\frac12\sum_{i,j}\|x_{ij}\|\,a_j
                           e^{|\Delta_{ij}|}\Delta_{ij}^2.
\]
Large centered log effects invalidate the assumption that the
first-order residual is close to the exact one, even if the
mean lesion incidence itself is small.

For a single voxel with intercept log mean $\eta_0$, two centered
covariate values $(-1,1)$, and coefficient $b$, let $a=e^{\eta_0}$.
The exact means are $(ae^{-b},ae^b)$; the approximation gives
$(a(1-b),a(1+b))$.
The intercept mean projection is $2a\cosh b$ versus $2a$.
The covariate mean projection is $2a\sinh b$ versus $2ab$.
Thus the intercept error begins at $ab^2$ and the covariate error
at $ab^3/3$. In the scalar basis case $B=1$, the correct
$K_a=a$, whereas squaring the weighted basis incorrectly gives $a^2$.
This example distinguishes both the intended approximation and the
unintended weighting change.

\subsection{An NB score approximation must also expand its denominator}
The manuscript's Taylor residual is a Poisson calculation.
For NB define the scalar score
$s(\eta)=(y-e^\eta)/(1+\alpha e^\eta)$.
At $\eta=\bar\eta_j$,
\[
 s(\bar\eta_j)=\frac{y-a_j}{1+\alpha a_j},\qquad
 s'(\bar\eta_j)=-\frac{a_j(1+\alpha y)}{(1+\alpha a_j)^2}.
\]
A first-order approximation to the actual NB score is therefore
\begin{equation}
 s_{ij}\approx
 \frac{Y_{ij}-a_j}{1+\alpha a_j}
 -\frac{a_j(1+\alpha Y_{ij})}{(1+\alpha a_j)^2}\Delta_{ij}.
 \label{lesion-nb-taylor}
\end{equation}
For a projected matrix expression set
$k_j=(1+\alpha a_j)^{-1}$ and
$c_j=a_j/(1+\alpha a_j)^2$. Then
\begin{align}
 S_{\rm lin}^{\rm NB}
 &=(\widetilde Y-\mathbf1_Ma^\top)\diag(k)
       -\Delta\diag(c)-\alpha(\widetilde Y\mathbin{\odot}\Delta)\diag(c),\\
 U_{\rm lin}^{\rm NB}&=\vr(Z^\top S_{\rm lin}^{\rm NB}B).
\end{align}
The final data-dependent term is absent for Poisson.
There is no automatic claim that it reduces to the same cached
products or computational cost as \eqref{lesion-taylor-score}.
Its validity is a local Taylor approximation with fixed $\alpha$,
not the exact NB update in (2.12).
As $\alpha\downarrow0$ it reduces to the Poisson expansion.

\section{Covariance for the spatially varying model}
\label{lesion-covariance}

\subsection{From a likelihood Taylor expansion to an inverse information}
For an unpenalized fit $\widehat\beta$, expand the score near its
population target $\beta_0$:
\[
 0=U(\widehat\beta)
 \approx U(\beta_0)-H(\beta_0)(\widehat\beta-\beta_0).
\]
Thus $\widehat\beta-\beta_0\approx H^{-1}U(\beta_0)$.
Under a correctly specified regular likelihood,
$\Var(U)=I$ and the large-sample sensitivity is $I$,
giving $\Cov(\widehat\beta)\approx I^{-1}$.
An observed-information alternative uses $H(\widehat\beta)^{-1}$.
These coincide algebraically for Poisson but not generally for NB.
Neither is the exact finite-sample covariance of a nonlinear estimator.

A fixed initialization preconditioner is not substituted for
$I(\widehat\beta)$ merely because it was used to optimize the fit.
Inference must use the curvature or sensitivity of the final
estimating equation, or explicitly quantify an approximation to it.

\subsection{Selected covariance by a solve}
For $C_\beta\in\R^{k\times RP}$, solve
\[
 IV=C_\beta^\top.
\]
Then $V=I^{-1}C_\beta^\top$ and
$C_\beta V=C_\beta I^{-1}C_\beta^\top$ is the desired
$k\times k$ covariance.
This needs $k$ right-hand sides rather than all $RP$ columns
of an inverse. A coordinate variance uses a coordinate row;
a voxelwise contrast uses the row derived in
Section~\ref{lesion-inference}.
Batch voxel contrasts to avoid storing a full $N\times N$
spatial covariance when only its diagonal is required.

\subsection{Design-supported block structure}
By \eqref{lesion-hessian-entry}, a spatial cross block is zero
if $Z_{ir}Z_{ir'}=0$ for every subject.
Connect columns of $Z$ when they co-occur in a row.
Different connected components give independent information blocks
after a coefficient permutation:
$H=\operatorname{blkdiag}(H_1,\ldots,H_K)$.
Direct multiplication proves
$H^{-1}=\operatorname{blkdiag}(H_1^{-1},\ldots,H_K^{-1})$
when each block is invertible.

An intercept column and continuously varying covariates usually
connect the SV maps, so group-style independent spatial blocks
are not generally available here.
No direct edge is weaker than no connecting path:
for $\bigl(\begin{smallmatrix}a&b&0\\b&c&e\\0&e&f\end{smallmatrix}\bigr)$
the $(1,3)$ inverse entry is $be/\det H$, generally nonzero.
Accidental zero entries at one fit should not define reusable
structural blocks.

\subsection{Joint dispersion adjustment}
If $\alpha$ is jointly estimated, write its observed information as
\[
 \mathcal I_{\rm obs}
 =\begin{pmatrix}H_{\beta\beta}&h\\h^\top&a_\alpha\end{pmatrix},
 \quad h=H_{\beta\alpha},\quad a_\alpha=H_{\alpha\alpha}.
\]
For a system with right-hand side $(b,0)$, the nuisance equation gives
$x_\alpha=-a_\alpha^{-1}h^\top x_\beta$.
Substitution into the regression equation gives
\[
 (H_{\beta\beta}-h a_\alpha^{-1}h^\top)x_\beta=b.
\]
Hence the regression block of the joint inverse is
\begin{equation}
 (\mathcal I_{\rm obs}^{-1})_{\beta\beta}
 =(H_{\beta\beta}-h a_\alpha^{-1}h^\top)^{-1}.
 \label{lesion-dispersion-schur}
\end{equation}
This requires the stated inverses; positive-definite joint curvature
gives a PSD information reduction and a nonnegative covariance
adjustment relative to treating dispersion as known.

For the correctly specified independent NB count likelihood,
$\E f_{\eta\alpha}=0$ because \eqref{lesion-nuisance-alpha}
is proportional to $Y-\mu$.
Regression and dispersion are therefore orthogonal in its
expected information, but the observed cross block need not vanish.
Alternating regression and dispersion updates does not itself imply
orthogonality or supply the joint adjustment.
For robust inference, stack both score equations and form their
joint sandwich; adjusting only the regression bread does not account
for nuisance-score variability.

\section{Subject-clustered sandwich: matching score, bread, and normalization}
\label{lesion-sandwich}

\subsection{Derivation with summed rather than averaged quantities}
Let $U_i$ be the chosen family's subject score and
$U=\sum_iU_i$. Define
\[
 A=-\sum_i\frac{\partial U_i}{\partial\beta^\top},\quad
 J=\sum_iU_iU_i^\top.
\]
At a suitable target, the score linearization gives
$\widehat\beta-\beta_0\approx A^{-1}\sum_iU_i$.
For independent subjects and valid moment assumptions, the empirical
covariance is
\begin{equation}\label{lesion-sandwich-formula}
 V_{\rm robust}=A^{-1}JA^{-\top}.
\end{equation}
No further $1/M$ multiplies this formula: $A$ and $J$ are
already sums. If $\bar A=A/M$ and $\bar J=J/M$,
$(1/M)\bar A^{-1}\bar J\bar A^{-\top}=A^{-1}JA^{-\top}$.
This checks the normalization explicitly.

For the likelihood score, $A=H$.
The NB expected-sensitivity weights in
\eqref{lesion-nb-fisher} are an asymptotically justified alternative
under the correct conditional mean and appropriate regularity.
They are not a license to change the score to the Poisson residual.

\subsection{The family-specific subject score}
The correct scores are
\[
 U_i^{\rm P}=X_i^\top(Y_i-\mu_i),\qquad
 U_i^{\rm NB}
 =X_i^\top\diag\{(1+\alpha\mu_i)^{-1}\}(Y_i-\mu_i).
\]
The manuscript's (2.18) uses the Poisson form for both families.
To interpret it as an NB likelihood sandwich, replace the score
by the NB expression above and use its matching sensitivity.
A common log link does not make the two likelihood scores identical.

Write $U_i=\sum_j u_{ij}$ and expand its outer product:
\[
 U_iU_i^\top
 =\sum_j u_{ij}u_{ij}^\top+
                  \sum_{j\ne k}u_{ij}u_{ik}^\top.
\]
Subject clustering retains the cross-voxel terms.
A cellwise sum of outer products drops them and is not robust
to arbitrary residual dependence within a lesion image.
If subjects themselves are dependent, the independent unit must
be enlarged appropriately; more voxels do not supply more
independent subjects.

Consistency also needs the relevant mean and law-of-large-numbers
conditions. Under a misspecified fixed design, uncentered score
outer products with nonzero individual score means cannot
automatically be equated with sampling variances.
A sandwich estimates variability around the estimator's target;
it does not remove mean misspecification or approximation bias.

\subsection{Contrast sandwich as a Gram matrix of projected scores}
Stack subject scores as columns,
$\mathcal U=[U_1,\ldots,U_M]\in\R^{RP\times M}$,
so $J=\mathcal U\mathcal U^\top$.
This is the column orientation of the matrix called $U$ in
manuscript (S2.48); the calligraphic symbol distinguishes it from
the total score $U=\sum_iU_i$.
For symmetric bread $A$, solve $AV=C_\beta^\top$ and let
$G=\mathcal U^\top V$.
Then
\[
 G^\top G
 =V^\top J V
 =C_\beta A^{-1}JA^{-1}C_\beta^\top.
\]
This obtains the contrast covariance without constructing a dense meat
or full inverse. Equivalently accumulate $g_i g_i^\top$ with
$g_i=C_\beta A^{-1}U_i$.
For a nonsymmetric estimating-equation bread solve
$A^\top V=C_\beta^\top$ instead.

The Gram representation is PSD and has rank at most $\min(M,k,d)$.
At an exactly solved unpenalized score equation,
$\sum_iU_i=0$, giving the stronger bound $\min(M-1,k,d)$.
A scalar degrees-of-freedom correction cannot increase that rank
or compensate for an inadequate number of independent subjects.

\subsection{Expand the manuscript's QR sandwich calculation}
For Poisson, define the paper's weighted subject design
$X_{w,i}=W_i^{1/2}X_i$, with $W_i=\diag(\mu_i)$.
The stacked weighted design $X_w$ satisfies
\[
 X_w^\top X_w=\sum_iX_{w,i}^\top X_{w,i}
            =\sum_iX_i^\top W_iX_i=A.
\]
This is (S2.44)--(S2.47). The supplement calls the bread $B$
and the meat $M$; we retain $A$ and $J$ to avoid confusing
them with the spatial basis and number of subjects.

For full-column-rank $X_w$, take a thin QR factorization
\[
 X_w=\mathcal Q\mathcal R,\qquad
 \mathcal Q^\top\mathcal Q=I_{RP},\qquad
 A=\mathcal R^\top\mathcal R .
\]
Here $\mathcal R$ is \emph{upper} triangular in the standard QR
convention, not lower triangular as described after (S2.48).
To compute $\mathcal S_U=A^{-1}\mathcal U$, solve
\[
 \mathcal R^\top T_U=\mathcal U,\qquad
 \mathcal R\mathcal S_U=T_U.
\]
Substitution gives
$\mathcal R^\top\mathcal R\mathcal S_U=\mathcal U$,
which verifies the solve order. The supplement names these
intermediates $Z$ and $S$; $T_U,\mathcal S_U$ keep them separate
from the study design and contrast count here.
Now multiply:
\begin{align}
 A^{-1}JA^{-1}
 &=A^{-1}\mathcal U\mathcal U^\top A^{-1}\\
 &=(A^{-1}\mathcal U)(A^{-1}\mathcal U)^\top
   =\mathcal S_U\mathcal S_U^\top.
 \label{lesion-qr-sandwich}
\end{align}
This derives (S2.49)--(S2.50) without constructing a dense meat.
For only $k$ contrasts, the smaller solves in the previous
subsection avoid storing all $RP\times M$ solved score columns.

For NB expected-sensitivity bread the same algebra uses
$W_i=\diag\{\mu_i/(1+\alpha\mu_i)\}$ and the NB subject scores.
For observed NB bread, use the observed weights instead.
Changing the bread's weights does not eliminate the NB residual
denominator in $\mathcal U$.

The full $MN\times RP$ stack need not be stored.
If $\mathcal R_{\rm old}$ represents already processed weighted
rows, QR-factor the stack of $\mathcal R_{\rm old}$ and the next
block $X_{w,\rm next}$. The resulting factor obeys
\[
 \mathcal R_{\rm new}^\top\mathcal R_{\rm new}
 =\mathcal R_{\rm old}^\top\mathcal R_{\rm old}
                   +X_{w,\rm next}^\top X_{w,\rm next}.
\]
This proves the blockwise accumulation of the same bread factor.
With a thin SVD $X_w=Q_wD_wV_w^\top$, the equivalent solve is
$A^{-1}\mathcal U=V_wD_w^{-2}V_w^\top\mathcal U$ when every
singular value is positive.
Since $\kappa_2(A)=\kappa_2(X_w)^2$, weighted-design QR/SVD can
avoid the extra rounding error of explicitly forming a Gram
matrix, but cannot create information in a singular design.

\section{Voxelwise inference under the manuscript's spatial-effect setup}
\label{lesion-inference}

\subsection{Linear map contrasts and dimension checks}
Let $C\in\R^{S\times R}$ specify the tested spatial covariate effects.
The map vector at voxel $j$ is
the paper's $\varphi_j=(I_R\otimes B_j^\top)\beta$.
Across voxels, retain its contrast-map notation
$\bar\eta_C=(C\otimes B)\beta$ and
$\bar\mu_C=\exp(\bar\eta_C)$.
The bar with subscript $C$ labels these contrasts; it does not
mean the subject averaging used to define $\bar\eta$ in the
Taylor expansion.
Multiplying by $C$ gives
\[
 T_j=C\otimes B_j^\top\in\R^{S\times RP},\quad
 \widehat\delta_j=T_j\widehat\beta-\delta_{0j},\quad
 V_j=T_j V_\beta T_j^\top.
\]
This supplies the dimensions for manuscript (2.15)--(2.16).
For $S$ independent estimable constraints and nonsingular consistent
$V_j$, $Q_j=\widehat\delta_j^\top V_j^{-1}\widehat\delta_j$
has the usual asymptotic $\chi_S^2$ null law.
To derive the degrees of freedom, factor $V_j=L_jL_j^\top$.
Under the null, $L_j^{-1}\widehat\delta_j$ is asymptotically a
vector of $S$ independent standard normals, and $Q_j$ is its
squared norm.
For one contrast, preserve the paper's signed statistic:
\[
 W_j=\frac{(C\otimes B_j^\top)\widehat\beta-\delta_{0j}}
 {\sqrt{(C\otimes B_j^\top)V_\beta(C\otimes B_j^\top)^\top}},
 \qquad Q_j=W_j^2 .
\]
The manuscript's (2.15)--(2.16) use $\delta_{0j}=0$.
The signed $W_j$ is asymptotically $N(0,1)$, whereas its square
is asymptotically $\chi_1^2$; these names are not interchanged here.
Reduce redundant constraints explicitly rather than using a
pseudoinverse to hide an unidentified test.
Contrast rows need not sum to one: differences have row sum zero.

\subsection{The vectorized Wald map uses covariance diagonals}
For $S=1$ and a zero null, set $T=C\otimes B\in\R^{N\times RP}$.
Then $\widehat{\bar\eta}_C=T\widehat\beta$ and its covariance
estimate is $TV_\beta T^\top$, an $N\times N$ matrix.
Only its diagonal contains the marginal variances required for
the paper's vector of signed statistics:
\[
 \mathbf W=
 \widehat{\bar\eta}_C\ \mathbin{/}\
 \sqrt{\diag(TV_\beta T^\top)}.
\]
Here $\diag$ extracts the diagonal of a square matrix, and
division and square root are elementwise. The full covariance
matrix under a square root is not an $N$-vector of standard
errors; this makes that step in (2.16) dimensionally explicit.
For a two-sided normal reference, $P_j=2\{1-\Phi(|W_j|)\}$,
where $\Phi$ is the standard normal distribution function.
The diagonal can be computed by the selected solves above
without forming the full map covariance.

\subsection{Comparing independently fitted groups}
The manuscript's group comparison can fit the SV model independently
in disjoint groups with coefficient vectors $\beta_g,\beta_h$.
If their data and fitting decisions are genuinely independent,
their joint covariance is $\operatorname{blkdiag}(V_g,V_h)$.
For a fixed covariate-effect contrast $c\in\R^R$,
\[
 T_j=[\,c^\top\otimes B_j^\top\quad
                  -c^\top\otimes B_j^\top\,],\qquad
 V_{\delta j}=(c^\top\otimes B_j^\top)(V_g+V_h)(c\otimes B_j).
\]
Equivalently, stack $(\beta_g^\top,\beta_h^\top)^\top$ and use
$C_{\rm pair}=(c^\top,-c^\top)\in\R^{1\times2R}$.
This is a row contrast; a terminal transpose would instead make
it a $2R\times1$ column.
Unequal group sizes enter $V_g,V_h$; equality of a scientific
effect under a null is not a claim that finite-sample estimates or
their variances are independent of sample size.
If groups share subjects, fitted nuisance coefficients, or a coupled
selection procedure, zero cross-covariance needs reconsideration.
This separately fitted SV comparison is not the shared-global
multi-group CBMR parameterization of the other paper.

\subsection{Polynomial covariates and nonlinear intensity summaries}
If an age covariate $t$ enters through $t,t^2,t^3$, label those
three components of the voxel coefficient vector as
$\varphi_{j1},\varphi_{j2},\varphi_{j3}$.
Their log effect is
$\varphi_{j1}t+\varphi_{j2}t^2+\varphi_{j3}t^3$ and its local derivative is
\[
 \frac{\partial\eta_j(t)}{\partial t}
 =\varphi_{j1}+2t\varphi_{j2}+3t^2\varphi_{j3}.
\]
This is a linear coefficient contrast at any specified $t$,
with covariate row $(1,2t,3t^2)$ in the polynomial coordinates.
A finite change from $t_0$ to $t_1$ instead uses
$(t_1-t_0,t_1^2-t_0^2,t_1^3-t_0^3)$.
Its intensity ratio is the exponential of that log contrast;
it is not generally the exponential of a single linear-age coefficient.

For a specified subject design row $Z_*^\top$, put
$x_{*j}=Z_*\otimes B_j$.
Then $\mu_{*j}=e^{x_{*j}^\top\beta}$ has
Jacobian $\partial\mu_{*j}/\partial\beta^\top=\mu_{*j}x_{*j}^\top$.
Therefore
$\Cov(\widehat\mu_{*j},\widehat\mu_{*k})
\approx\mu_{*j}\mu_{*k}x_{*j}^\top V_\beta x_{*k}$.
For a scalar log contrast $\delta$, the delta-method variance
of $e^{\widehat\delta}$ is approximately
$e^{2\delta}\Var(\widehat\delta)$.
These Jacobians are exact; the covariance transformations are
first-order approximations.
An intensity ratio is a risk ratio only when the fitted intensities
can validly be interpreted as the conditional binary means.
For the explicitly conditioned logistic model,
$e^\delta$ is instead an odds ratio.

\subsection{Multiplicity and bootstrap remain separate statistical steps}
Pointwise Wald $p$-values do not automatically control error over
many voxels. Benjamini--Hochberg targets FDR under its required
dependence conditions, not general FWER control.
Neither faster Hessians nor a different spatial basis proves those
conditions or calibrates a finite-sample tail.

A nonparametric bootstrap should resample independent subjects with
their whole binary images and covariates, or the appropriate larger
independent clusters. A fixed-design parametric bootstrap keeps
covariates fixed and simulates the specified response model.
For a null test, fit and simulate the null rather than the unrestricted
alternative, and repeat the estimation, dispersion fitting, and
data-driven tuning steps that define the reported estimator.
If Poisson/NB counts are simulated, the experiment assesses that
count-model bootstrap; it does not silently become a Bernoulli
lesion bootstrap. A spatial binary simulator must preserve the
intended within-subject dependence.
Failed refits must be reported and diagnosed, not silently removed
from the reference distribution.

\section{Supplement S2.3: a complete toy slope-map variance derivation}
\label{lesion-toy-variance}

\subsection{Translate the supplement's local symbols and coefficient order}
Only in this section, follow S2.3 with \emph{scalar} subject values
$Z_i$ and spatial values $B_j$, and two-component vectors
\[
 x_i=(1,Z_i)^\top,\qquad b_j=(1,B_j)^\top,\qquad
 \beta_{\rm toy}=(\alpha,\gamma,\delta,\theta)^\top.
\]
These scalar $Z_i,B_j$ are not the main text's design columns.
Likewise, the supplement's $\alpha$ is now a log-intercept, not
NB dispersion; $\gamma,\delta,\theta$ are scalar toy coefficients,
not the initialization vector or log-dispersion coordinate.
We retain these local names so each step can be compared with S2.3.
Let $\mathsf X$ and $\mathsf B$ have rows $x_i^\top$ and $b_j^\top$;
these are the local matrices called $X$ and $B$ in that supplement.

Expand the supplement's basis-first Kronecker product:
\[
 b_j\otimes x_i=(1,Z_i,B_j,Z_iB_j)^\top,\qquad
 \log\mu_{ij}=\alpha+\gamma Z_i+\delta B_j+\theta Z_iB_j.
\]
The main text instead orders $x_i\otimes b_j=(1,B_j,Z_i,Z_iB_j)^\top$.
Thus its map-major coefficient vector is
$(\alpha,\delta,\gamma,\theta)^\top$, obtained by interchanging
coordinates two and three.
If $\Pi$ is that permutation matrix, then
$\beta_{\rm main}=\Pi\beta_{\rm toy}$ and
$I_{\rm main}=\Pi I_{\rm toy}\Pi^\top$.
Contrasts must be permuted with the coefficients; the variance
of the same scientific slope is then unchanged.

\subsection{Separate the subject and spatial Gram matrices}
Assume $\sum_i Z_i=0$, as in the supplement, and define
\[
 S_Z=\sum_i Z_i^2,\qquad S_B=\sum_j B_j^2.
\]
Direct multiplication, without mixing the two sets of indices, gives
\[
 \mathsf X^\top\mathsf X=
 \begin{pmatrix}M&0\\0&S_Z\end{pmatrix},\qquad
 \mathsf B^\top\mathsf B=
 \begin{pmatrix}N&\sum_jB_j\\\sum_jB_j&S_B\end{pmatrix}.
\]
Subject centering eliminates $\sum_iZ_i$ from the first matrix.
It does not replace that zero by a spatial sum $\sum_jB_j$,
and does not replace the subject count $M$ by the voxel count $N$.

The exact Poisson information in the supplement's order is
\[
 I_{\rm toy}=\sum_{i,j}\mu_{ij}
                     (b_jb_j^\top)\otimes(x_ix_i^\top).
\]
If the subject-dependent log effects $(\gamma+\theta B_j)Z_i$
are small, approximate
$\mu_{ij}\approx\mu_j=\exp(\alpha+\delta B_j)$.
This equality is exact at $\gamma=\theta=0$.
Pulling the separate sums out of the Kronecker product gives
\[
 I_{\rm toy}\approx
 \left(\sum_j\mu_j b_jb_j^\top\right)\otimes
 \left(\sum_i x_ix_i^\top\right).
\]
Write the supplement's scalar moments $A,B,C$ as
\[
 A_\mu=\sum_j\mu_j,\quad
 B_\mu=\sum_jB_j\mu_j,\quad
 C_\mu=\sum_jB_j^2\mu_j,
\]
to avoid collisions with the spatial design and contrast matrix.
With $K_\mu=\bigl(\begin{smallmatrix}A_\mu&B_\mu\\
B_\mu&C_\mu\end{smallmatrix}\bigr)$ and
$D_Z=\diag(M,S_Z)$, the four-by-four information is
\begin{equation}
 I_{\rm toy}\approx K_\mu\otimes D_Z
 =\begin{pmatrix}
 MA_\mu&0&MB_\mu&0\\
 0&S_ZA_\mu&0&S_ZB_\mu\\
 MB_\mu&0&MC_\mu&0\\
 0&S_ZB_\mu&0&S_ZC_\mu
 \end{pmatrix}.
 \label{lesion-toy-information}
\end{equation}

\subsection{Invert the two factors and project to the slope}
The slope at voxel $j$ is the supplement's
$m_j=\gamma+B_j\theta=(b_j\otimes c)^\top\beta_{\rm toy}$,
where $c=(0,1)^\top$.
Assume $S_Z>0$ and
$\Delta_\mu=A_\mu C_\mu-B_\mu^2>0$.
Then
\[
 K_\mu^{-1}=\frac1{\Delta_\mu}
       \begin{pmatrix}C_\mu&-B_\mu\\-B_\mu&A_\mu\end{pmatrix},
 \qquad D_Z^{-1}=\diag(1/M,1/S_Z).
\]
The mixed-product identity gives
\begin{align}
 \Var(\widehat m_j)
 &\approx(b_j\otimes c)^\top
             (K_\mu^{-1}\otimes D_Z^{-1})(b_j\otimes c)\\
 &=(b_j^\top K_\mu^{-1}b_j)(c^\top D_Z^{-1}c)\\
 &=\frac{C_\mu-2B_\mu B_j+A_\mu B_j^2}
                {S_Z(A_\mu C_\mu-B_\mu^2)}.
 \label{lesion-toy-variance-formula}
\end{align}
The approximation combines the small-slope information
factorization with inverse-information asymptotics for the estimator.
The matrix inversion and contrast projection are exact algebra.

For a more interpretable form, define
\[
 \bar B_\mu=B_\mu/A_\mu,\qquad
 S_{BB,\mu}=\sum_j\mu_j(B_j-\bar B_\mu)^2
           =C_\mu-B_\mu^2/A_\mu.
\]
Expanding and completing the square in the numerator yields
\[
 \Var(\widehat m_j)\approx\frac1{S_Z}
       \left\{\frac1{A_\mu}
                 +\frac{(B_j-\bar B_\mu)^2}{S_{BB,\mu}}\right\}.
\]
The determinant condition is $A_\mu S_{BB,\mu}>0$:
the spatial feature must vary on positive-weight voxels.
Neither a constant spatial feature nor a constant subject covariate
can be repaired by an inverse formula.

\subsection{Constant means, pooling, and a correction to the printed formula}
When $\mu_j\approx\bar\mu$, let
$\bar B=N^{-1}\sum_jB_j$ and
$S_{BB}=\sum_j(B_j-\bar B)^2$.
Then $A_\mu=N\bar\mu$, $\bar B_\mu=\bar B$, and
$S_{BB,\mu}=\bar\mu S_{BB}$, giving
\begin{equation}
 \Var(\widehat m_j)\approx
 \frac1{\bar\mu S_Z}
 \left\{\frac1N+\frac{(B_j-\bar B)^2}{S_{BB}}\right\}.
 \label{lesion-toy-constant}
\end{equation}
The extra factor involving
$NS_Z-(\sum_jB_j)^2$ in the supplied S2.3 derivation does not
follow from its centered subject design. It mixes a voxel moment
into the subject Gram matrix. In particular the final denominator
must not replace $S_Z$ by $S_Z-N\bar B^2$.
Equations \eqref{lesion-toy-information}--\eqref{lesion-toy-constant}
show where the two dimensions remain separate.

For a hand check, take $M=3$, $Z_i=(-1,0,1)$, $N=4$,
$B_j=(0,1,2,3)$, and $\bar\mu=0.1$.
Then $S_Z=2$, $\bar B=1.5$, and $S_{BB}=5$.
The corrected inverse-information slope variances are
$(3.5,1.5,1.5,3.5)$; the erroneous factor
$S_Z-N\bar B^2=-7$ would give negative values.
These are algebraic audit values, not a claim of accurate
small-sample inference with only three subjects.

Under the same independent Poisson assumptions, separate
voxelwise intercept/slope fits have inverse-information slope
variance $1/(\bar\mu S_Z)$.
The spatial result multiplies it by
$h_j=1/N+(B_j-\bar B)^2/S_{BB}$.
These are the diagonal leverages of the two-column spatial
design: $0\le h_j\le1$ and $\sum_jh_j=2$.
This quantifies pooling when the spatial slope restriction is
appropriate; it is not a guarantee under a misspecified slope
map or residual spatial dependence.
For binary or correlated lesion images, the applicable
score-variability calculation remains the subject sandwich,
not an automatic identification of Poisson information with variance.

\section{The mass-univariate logistic and Firth comparator}
\label{lesion-logistic}

\subsection{Logistic score and Hessian at one voxel}
For the comparator in manuscript S1.1, fit each voxel separately.
Its symbols $(X,x_i,\beta_j,\eta_{ij})$ correspond here to
$(Z,Z_i,\varphi_j,\zeta_{ij})$, respectively.
This preserves the main text's design and voxel-coefficient names
while reserving $X,\beta$ for the full spatial model.
The comparator's $\varphi_j$ is fitted with a logistic link here;
it is not asserted to equal the Poisson MUM estimate in (2.5).
Use coefficient $\varphi_j\in\R^R$, predictor
$\zeta_{ij}=Z_i^\top\varphi_j$, and
$p_{ij}=e^{\zeta_{ij}}/(1+e^{\zeta_{ij}})$.
The Bernoulli log-likelihood is
\begin{align}
 \ell_j
 &=\sum_i\{Y_{ij}\log p_{ij}+(1-Y_{ij})\log(1-p_{ij})\}\\
 &=\sum_i\{Y_{ij}\zeta_{ij}-\log(1+e^{\zeta_{ij}})\}.
\end{align}
Since $\partial_\zeta p=p(1-p)$,
\[
 U_j=Z^\top(Y_{\cdot j}-p_{\cdot j}),\quad
 H_j=I_j=Z^\top\diag\{p_{\cdot j}(1-p_{\cdot j})\}Z.
\]
This $R\times R$ matrix is for one voxel's comparator, not the
$RP\times RP$ spatial model. Restricting the logistic coefficient
field to a basis would instead use the full spatial design and
weights $p_{ij}(1-p_{ij})$.
Neither formulation has the Poisson weight $\mu_{ij}$.

\subsection{Firth's determinant term is not a fixed quadratic penalty}
The manuscript's Firth comparator maximizes
$\ell_{F,j}=\ell_j+\tfrac12\log\det I_j$.
Assume $I_j$ is nonsingular. The identity
$d\log\det I=\tr(I^{-1}dI)$ gives
\begin{equation}
 (U_{F,j})_a=(U_j)_a+\tfrac12\tr(I_j^{-1}I_{j,a}),
 \qquad I_{j,a}=\partial_{\varphi_{j,a}}I_j .
\end{equation}
Write $w_i=p_i(1-p_i)$, suppressing the fixed voxel $j$.
Then
\begin{align}
 \partial_{\zeta_i}w_i&=w_i(1-2p_i),\\
 \partial^2_{\zeta_i}w_i
 &=w_i(1-2p_i)^2-2w_i^2,\\
 I_{j,a}
 &=\sum_i w_i(1-2p_i)Z_{ia}Z_iZ_i^\top,\\
 I_{j,ab}
 &=\sum_i\{w_i(1-2p_i)^2-2w_i^2\}
                         Z_{ia}Z_{ib}Z_iZ_i^\top .
\end{align}
The first derivative also has a directly interpretable
leverage form. Define $h_i=w_i Z_i^\top I_j^{-1}Z_i$.
Using $\tr(I_j^{-1}Z_iZ_i^\top)=Z_i^\top I_j^{-1}Z_i$,
\begin{align}
 \tfrac12\tr(I_j^{-1}I_{j,a})
 &=\tfrac12\sum_i w_i(1-2p_i)Z_{ia}
                               Z_i^\top I_j^{-1}Z_i\\
 &=\sum_i h_i(1/2-p_i)Z_{ia},\\
 U_{F,j}&=Z^\top\{Y_{\cdot j}-p_{\cdot j}
                         +h\mathbin{\odot}(1/2-p_{\cdot j})\}.
\end{align}
Here $Y_{\cdot j}=(Y_{1j},\ldots,Y_{Mj})^\top$ is the vector
at voxel $j$, not the initialization's summed margin $(Y_B)_j$.
For a full-rank design,
$\sum_i h_i=\tr(I_j^{-1}I_j)=R$.
This expands the determinant adjustment, rather than treating it
as a constant ridge term.
Differentiating $I^{-1}I=I_R$ gives
$\partial_b I^{-1}=-I^{-1}I_bI^{-1}$.
Therefore the full observed information of the penalized comparator is
\begin{equation}
 (H_{F,j})_{ab}=(H_j)_{ab}
 -\tfrac12\tr\{I_j^{-1}I_{j,ab}
                   -I_j^{-1}I_{j,b}I_j^{-1}I_{j,a}\}.
\end{equation}
The determinant term depends on the coefficients and has its own
curvature. It is not the fixed roughness penalty from the
published coordinate-based CBMR model.
The sampling covariance of a bias-reduced estimator also requires
its inferential assumptions; an inverse penalized Hessian is
not a universal finite-sample standard-error formula.

\section{Numerical conditioning and explicitly optional regularization}

\subsection{Symmetric eigenvalues replace an unnecessary general SVD}
For symmetric nonsingular $H=Q\Lambda Q^\top$,
$H^\top H=Q\Lambda^2Q^\top$, so its singular values are
$|\lambda_a|$. Thus
\[
 \kappa_2(H)=\frac{\max_a|\lambda_a|}{\min_a|\lambda_a|}.
\]
For SPD information the absolute values are unnecessary.
An indefinite matrix may have a finite condition number but not a
valid inverse-covariance interpretation.
For disconnected blocks take the largest absolute eigenvalue over
all blocks divided by the smallest over all blocks, not the
largest within-block condition number.
For example, blocks $I$ and $10^6I$ each have condition number one,
while their combined matrix has condition number $10^6$.

For exact $P_0=K_Z\otimes K_B$,
$\kappa_2(P_0)=\kappa_2(K_Z)\kappa_2(K_B)$.
This diagnostic describes that preconditioner; it need not describe
the final nonseparable information.
Solving with a factorization is preferable to explicitly forming
an inverse when only directions or contrasts are requested.
Symmetrizing roundoff does not repair genuine singularity.

\subsection{A fixed quadratic penalty would be an additional model choice}
The manuscript's core spatial model reduces dimension using a fixed
basis; it does not thereby acquire the other paper's roughness penalty.
If one deliberately adds $\mathcal P(\beta)=\tfrac12\beta^\top K\beta$
with fixed symmetric $K\succeq0$, then
\[
 \nabla(-\ell+\mathcal P)=-U+K\beta,\qquad
 H_{\rm pen}=H+K.
\]
Linearizing $U(\widehat\beta)-K\widehat\beta=0$ gives
\[
 (H+K)(\widehat\beta-\beta_0)\approx U(\beta_0)-K\beta_0.
\]
Its local sampling covariance is
$(H+K)^{-1}J(H+K)^{-\top}$, not automatically $(H+K)^{-1}$,
and its local bias includes $-(I+K)^{-1}K\beta_0$.
Under a correctly specified count likelihood,
\[
 (I+K)^{-1}-(I+K)^{-1}I(I+K)^{-1}
 =(I+K)^{-1}K(I+K)^{-1}.
\]
Inverse penalized curvature is instead the usual local Gaussian
posterior curvature when the penalty is interpreted as a prior.
Neither expression includes data-selected penalty uncertainty.

Adding a numerical ridge only when computing a covariance changes
the inverse without changing the estimator's fitted equation.
It cannot be described as new data identification.
For PSD $H,K$, $H+K\succ0$ exactly when their null spaces share
no nonzero direction; a roughness penalty may leave a constant
confounding direction untouched.

\subsection{Computational accounting for the manuscript's alternatives}
The following orders use dense arithmetic and a fixed basis.
They separate fitting calculations from the covariance output;
they are not measured timings.
\begin{center}\small
\begin{tabularx}{\textwidth}{TT}
\toprule
Calculation & Leading operation count and storage implication\\
\midrule
Exact batched Poisson mean and score &
$O(MNP+MRP)$ work with subject batching; no explicit
$MN\times RP$ design is needed\\
Exact independent-cell information &
$O(MNP^2+MR^2P^2)$ for studywise basis products and block assembly;
dense output still has $(RP)^2$ entries\\
Fixed separable Poisson preconditioner &
Factor work $O(P^3+R^3)$ after forming its two weighted Gram matrices;
factor storage $O(P^2+R^2)$\\
Poisson Taylor score after fixed projections &
$O(NP^2+R^2P+RP^2)$ per iterate for $K_a$, matrix products and
the new expansion point; no full current $M\times N$ mean array\\
Dense final information factorization &
$O((RP)^3)$ work and $O((RP)^2)$ storage unless additional exact
structure is available\\
$k$ contrast covariance solves &
$O(k(RP)^2)$ after a dense factorization, with $RP\times k$
solve output instead of a full inverse\\
Exact Hessian--vector product &
$O(MRP+MNP)$ per product with current weights; no dense information
matrix, but iterative solves need convergence control\\
\bottomrule
\end{tabularx}
\end{center}
The exact-score order follows by evaluating
$\widetilde\eta=(Z\widetilde\beta)B^\top$ and
$Z^\top(\widetilde Y-\widetilde\mu)B$, choosing that multiplication order.
The information order follows from one $O(NP^2)$ weighted basis
product per subject and its distribution into $R^2$ coefficient blocks.
For the Taylor row, cache $Z^\top\widetilde YB$, $Z^\top \widetilde Z$,
and $Z^\top\mathbf1_M$ once.
At an iterate form $a$ and $K_a=B^\top\diag(a)B$, then multiply
$(Z^\top \widetilde Z)\widetilde\beta K_a$ using matrices of sizes $R$ and $P$.
These costs do not apply unchanged to the data-dependent
NB correction in \eqref{lesion-nb-taylor}.
Input storage, sparse count projections, basis sparsity, and batching
affect actual resource usage, and a small approximate score cost does
not establish its accuracy.

\section{Hand calculations and a paper-specific review checklist}

\subsection{A two-subject, two-voxel spatial model}
Take $B=I_2$ and
$Z=\bigl(\begin{smallmatrix}1&0\\1&1\end{smallmatrix}\bigr)$.
Order $\beta=(\beta_{1,1},\beta_{1,2},\beta_{2,1},\beta_{2,2})^\top$,
with the first covariate the intercept.
At $\beta=0$, all $\mu_{ij}=1$ and
\[
 X=\begin{pmatrix}1&0&0&0\\0&1&0&0\\1&0&1&0\\0&1&0&1\end{pmatrix},
 \quad
 H_{\rm P}=
 \begin{pmatrix}2&0&1&0\\0&2&0&1\\1&0&1&0\\0&1&0&1\end{pmatrix}.
\]
This is an algebraic evaluation point, not a recommended fitted
binary mean. With binary counts $(0,1,1,0)$ and $\alpha=1$,
the NB observed weights are $(1,2,2,1)/4$, yielding
\[
 H_{\rm NB}=
 \begin{pmatrix}
 3/4&0&1/2&0\\0&3/4&0&1/4\\1/2&0&1/2&0\\0&1/4&0&1/4
 \end{pmatrix},
 \qquad I_{\rm NB}=\tfrac12H_{\rm P}.
\]
Each entry is the sum of four weighted design-row products.
This makes the observed/expected distinction directly checkable
without a numerical differentiation package.

\subsection{Decisions that must be explicit when reading the manuscript}
\begin{center}\small
\begin{tabularx}{\textwidth}{lT}
\toprule
Location & Required interpretation or correction\\
\midrule
S1.2 versus (2.8) & Conditioning Poisson on zero/one gives logistic;
the unconditioned Poisson objective is different\\
(2.9)--(2.10) & NB is a count working model on binary observations,
not a literal positive-dispersion Bernoulli family\\
(2.11), Algorithm 1 & Identify the additive intercepts and embed the
$P$-vector fit into $PR$ coordinates explicitly\\
(2.12), Algorithm 1 & NB score has the denominator $1+\alpha\mu$\\
(2.13) & Exact for separable Poisson weights; approximate for general NB\\
Fixed-step iteration & SPD preconditioning alone does not guarantee
convergence of an undamped unit-step update\\
(2.14), S2.1 & Use centered $Q_MZ$, unweighted $B$ in the data
projection, vector $\beta$ in the Kronecker term, and preserve
the single-weight $B^\top\widetilde B$ product\\
(2.14), moving expansion point & Its approximate score can have a
nonsymmetric Jacobian; it is not automatically a likelihood gradient\\
(2.15)--(2.16) & Contrasts need not sum to one; ranks and dimensions
determine degrees of freedom. Keep $W_j$ signed and use marginal
covariance diagonals for the vectorized standard errors\\
(2.18) & Use the family's own subject score and matching bread\\
S2.2 & QR's $\mathcal R$ is upper triangular; retain the column-stacked
scores and correct two-solve order\\
S2.3 & Translate the basis-first coefficient order, keep the subject
Gram factor $\diag(M,S_Z)$, and do not insert voxel moments into it\\
\bottomrule
\end{tabularx}
\end{center}
These qualifications distinguish mathematical changes from performance
improvements. They do not claim to assess every simulation, empirical
conclusion, or implementation associated with the manuscript.
The count-model symbolic identities in \texttt{proofs/appendix.py}
and related repository modules support parts of this exposition;
the new paper-specific text and corrections are not automatically
covered by the existing Lean inventory.

\subsection{Editable source and independent build}
This document is self-contained: it does not require the combined
appendix or the other paper's appendix to be read.
Its editable source is \texttt{docs/cbmr\_paper\_appendix.tex}.
From the repository root:
\begin{verbatim}
make cbmr-paper-appendix
\end{verbatim}
Build both paper-specific appendices with \texttt{make paper-appendices}.
The original papers and combined expanded appendix are preserved.
\end{document}
